\begin{abstract}
We study problems related to connecting multi-interface networks of wireless devices. These problems can be modeled using graphs, where vertices represent the devices and edges represent potential communication  links. Each vertex can activate multiple interfaces, and a connection between two vertices is established if they share at least one common active interface.
However activating an interface induces a cost that depends on the type of the interface and on the vertex that activates it. We consider two problems arising in multi-interface networks: \coverage and \connectivity. In the \coverage problem, every connection defined in the network must be established, while in the \connectivity problem, it is only required that established connections form a connected spanning subgraph of the network. The solution should also minimize either the maximum cost incurred by a vertex or the total cost incurred by all vertices. In this work we are interested in approximating the former of the two cost criterions.

We model both problems using Integer Linear Programming (ILP) and we design approximation algorithms based on a randomized rounding of the solution of the linear programming relaxation. For the \coverage problem, this yields an $O(\log n)$-approximation algorithm, where $n$ is the number of vertices. This result is tight, since the problem generalizes \setcover. This improves upon the $O(b\cdot\log n)$-approximation algorithm for the min-sum version of the problem, where $b$ is a certain graph parameter which can be as large as $\Omega(n)$ [Algorithmica '12]. The same relaxation can also be used to get a $k$-approximation algorithm, where $k$ is the number of different interfaces. This generalizes a similar result for the homogeneous cost case, where the cost of an interface is the same for all vertices.
Our main result is an $O(\log^2 n)$-approximation algorithm for \connectivity, which is the first non-trivial approximation for this problem. The algorithm is based on a similar LP relaxation with additional cut constraints to ensure connectivity. The rounding procedure resembles the one for \coverage but requires a more careful analysis to ensure that the connectivity constraints are satisfied.
\end{abstract}

